\section{SI-Mapper: An Agentic Approach to Building Semantic Modeling}
\label{sec:implementation}

The previous section established that manual BACnet point mapping is costly, labor-intensive,
and resistant to automation with existing tools.
This section introduces the proposed solution --- agentic AI for automated HVAC semantic
modeling --- and traces its full development from initial concept through nine experimental
iterations to the final production architecture.

\subsection{Overview of the Approach}
\label{subsec:si-mapper-overview}

SI-Mapper takes two raw building data inputs, HVAC engineering drawings (PNG or PDF) and BACnet
point exports (CSV files), and produces two outputs: an ASHRAE 223P semantic model in Turtle (TTL)
format, and an interactive graph visualization of that model in a web interface.

The entire pipeline is orchestrated by an AI agent application guided by domain-specific skill instructions.
Figure~\ref{fig:pipeline} illustrates the data flow.

\begin{figure}[htbp]
\centering
\includegraphics[width=\textwidth]{pipeline.png}
\caption{SI-Mapper data pipeline --- from raw building data to semantic model and graph visualization.}
\label{fig:pipeline}
\end{figure}

The key enabler is multimodal AI.
Recent large language models like Gemini 3.1 and Claude Sonnet 4.6 can analyze engineering
drawings directly, identifying equipment symbols, ductwork connections, and spatial relationships from
visual inputs.
HVAC schematic diagrams, with their regular symbol vocabulary and consistent layout conventions, are
more tractable for AI recognition than general architectural floor plans.
The insight from SI-Mapper's development is that standard models underperform on domain-specific
drawing recognition until guided by detailed, domain-specific instructions --- the \emph{skills}
architecture described in Section~\ref{subsec:skills}.

SI-Mapper's novelty lies in combining four capabilities that no prior system integrates:
\begin{enumerate}[leftmargin=*, itemsep=2pt]
  \item \textbf{Multimodal drawing recognition:} AI analysis of HVAC engineering drawings to extract
        equipment and ductwork topology.
  \item \textbf{BACnet CSV mapping:} automated assignment of BACnet point data to the equipment
        identified in the drawings.
  \item \textbf{ASHRAE 223P generation:} production of a formally valid ASHRAE 223P ontology
        including ConnectionPoints, Connections, mediums, and telemetry links.
  \item \textbf{Iterative self-correction:} the agent compares its generated model against validation
        rules, identifies errors, and corrects them autonomously --- looping until the model is valid.
\end{enumerate}

This approach replaces weeks of expert engineering effort with an automated pipeline that runs in
minutes, producing a machine-readable semantic model that unlocks software portability, advanced
analytics, and grid integration for the building.

% -----------------------------------------------------------------------
\subsection{Development Methodology: Nine Experimental Iterations}
% -----------------------------------------------------------------------

The development of SI-Mapper followed an iterative experimental methodology.
Over nine distinct iterations, the system architecture evolved from simple function-calling
scripts to a sophisticated skills-based agentic AI system. Each iteration was driven by
the failures and limitations of the previous one.

\begin{figure}[htbp]
  \centering
  \includegraphics[width=\textwidth]{experiment-progression.png}
  \caption{Evolution of SI-Mapper architecture across nine experimental iterations.}
  \label{fig:experiment-progression}
\end{figure}

\textbf{Iteration 1 --- LLMs as Function Providers (Gemini 1.5 Flash).}
The first approach used LLMs as providers of intelligence embedded inside conventional
Python functions: the model was called for specific tasks such as image recognition or
text generation, while all business logic remained hardcoded by the developer.
This approach was straightforward but fundamentally limited --- the LLM had no awareness
of the application context and could not make autonomous decisions. Every new capability
required writing new code around it. LangChain was used as the integration library at
this stage.

\textbf{Iteration 2 --- Workflow-Based Agents (Gemini 1.5 Flash).}
The second iteration introduced workflows to define agent behavior as a sequence of
ordered steps, making it possible to better exploit the LLM's capabilities. However,
the business logic was still hardcoded inside the workflow definitions, and the agent
remained unable to adapt to unexpected situations. LangGraph was evaluated at this stage
but was abandoned due to its high commercial licensing costs, which made it unsuitable
for production deployment.

\textbf{Iteration 3 --- ADK Workflows (Gemini 2.0 Flash).}
The third iteration migrated to Google's Agent Development Kit (ADK) while retaining
the workflow-based execution model. A local RAG system was introduced to provide agents
with on-demand access to documentation for HVAC equipment classification. ADK had been
very recently released at this point, with limited documentation and community examples,
meaning AI coding assistants could not reliably help with the library --- a concrete
limitation that slowed development velocity.

\textbf{Iteration 4 --- Agentic Master Agent (Gemini 2.5 Flash).}
The fourth iteration represented the first true agentic attempt: a single master agent
given full autonomy to plan and execute all tasks. This introduced genuine decision-making
capability but exposed a critical failure mode --- the prompts required to express the
full HVAC pipeline were too long and complex for the model to follow reliably. A simple
while loop was not sufficient to keep the agent on track. This iteration also introduced
the CopilotKit frontend, the Agent-to-UI protocol for real-time agent observability, and
MCP tools for external tool management.

\textbf{Iteration 5 --- Plan-Act-Review Architecture (Gemini 2.5 Flash).}
Inspired by established multi-agent design patterns, the fifth iteration decomposed
execution into three sequential sub-agents: Planner, Actor, and Reviewer. This improved
output quality but the agent consistently terminated tasks before completing the full
pipeline. Adding a loop agent on top of the PAR structure improved results further but
was still insufficient --- the agent could not reliably complete all stages in a single
run.

\textbf{Iteration 6 --- Master Agent with PAR Sub-Agent (Gemini 2.5 Flash).}
The sixth iteration added a master agent that could choose to act directly or delegate to
the PAR structure as a sub-agent. Both the master and the sub-agent were given full
autonomy. Despite this added flexibility, the PAR sub-agent still could not produce
results of sufficient quality. Frontend tooling was extended with tool call and state
variable tracking to provide better visibility into the agent's reasoning.

\textbf{Iteration 7 --- Master Agent with Hybrid Workflows (Gemini 2.5 Flash).}
The seventh iteration replaced the PAR agent with a hybrid approach: specialized image
recognition workflows for the drawing analysis stage, combined with an agentic planner
for task orchestration. This created an architecture that was neither fully workflow-based
nor fully agentic. The recognition quality was not good enough, and the architectural
hybrid added complexity without proportional benefit. The entire drawing sequence was
ultimately converted back to a workflow, but results remained insufficient.

\textbf{Iteration 8 --- Parallel Agents (Gemini 2.5 Flash, then Gemini 3.0 Flash).}
The eighth iteration introduced parallel execution: one specialized agent per equipment
type running simultaneously, replacing the sequential recognition agents. This
significantly improved processing speed. A model upgrade to Gemini 3.0 Flash delivered
a meaningful accuracy improvement. Gemini 3.0 Pro was also tested and produced higher
quality results but required considerably more time, making it impractical for routine
use.

\textbf{Iteration 9 --- Skills-Based Master Agent (Gemini 3.1 / Claude Sonnet 4.6).}
The ninth and final iteration converted all sub-agents into skills --- specialized
instruction sets executed inline within the master agent's context. This architectural
simplification was made possible by the dramatically improved instruction-following
capability of the latest models. The agent was also given the ability to visit the
frontend via headless browser, capture a screenshot of the live HVAC canvas, and compare
it against the source drawing for visual self-correction. The architecture became
dramatically simpler, easier to maintain, and produced better results than any previous
iteration.

The key insight from this progression is that advances in foundation model capability made
simpler architectures viable. The skills-based approach became possible only when models
like Gemini 3.1 and Claude Sonnet 4.6 reached sufficient instruction-following capability
to execute complex, multi-step domain tasks within a single agent context.

% -----------------------------------------------------------------------
\subsection{Final Architecture}
% -----------------------------------------------------------------------

\begin{figure}[htbp]
  \centering
  \includegraphics[width=\textwidth]{architecture.png}
  \caption{SI-Mapper final architecture --- skills-based single master agent.}
  \label{fig:architecture}
\end{figure}

The final architecture organizes SI-Mapper into three main layers that together implement
the full pipeline from raw HVAC drawings to a validated ASHRAE 223P semantic model stored
in a graph database.

\textbf{Frontend (CopilotKit + Next.js).}
The frontend serves as a real-time operations dashboard. A CopilotKit Agent HUD streams
the agent's thoughts and actions as they occur, giving operators full visibility into what
the agent is doing and why. A Graphivac CAD canvas renders the internal HVAC grid in
real time as the agent places ducts and equipment. A Sigma.js graph viewer displays the
final semantic model after it is loaded into Neo4j. A code viewer shows snapshots of the
generated ASHRAE 223P Python code and the resulting TTL file at each stage of validation.

\textbf{Master Agent (Google ADK).}
The core of the system is a single \textit{LlmAgent} built on Google's Agent Development
Kit, running on Gemini 3.1 or Claude Sonnet 4.6. The agent has access to approximately
25 tools and six skills. It maintains full context across the entire pipeline --- from
drawing recognition through grid placement, BACnet metadata mapping, and ontology
generation and validation --- without handing off control to any sub-agent.

\textbf{External Services.}
Two external services underpin the system. Neo4j with the Neosemantics (n10s) plugin
stores the ASHRAE 223P ontology as a native RDF graph and provides Cypher query access
for graph traversal. The Graphivac REST API manages the CAD canvas, allowing the agent
to read and write the visual grid state bidirectionally.

\medskip
\noindent
The choice of a single-agent architecture over the multi-agent sub-agent approaches of
earlier iterations was driven by several compounding advantages.
A single agent avoids the coordination overhead inherent in multi-agent systems: there is
no inter-agent messaging, no shared state synchronization, and no risk of a sub-agent
misunderstanding its scope and terminating early.
Skills provide domain expertise inline, making the agent's behavior fully observable in a
single context window without the opacity of delegated sub-agent calls.
The agent naturally maintains full context across all pipeline stages, which is critical
for operations such as visual self-correction --- where the agent must compare a live
screenshot of the canvas against the source drawing and determine what to fix.
Finally, the architecture is dramatically simpler to maintain, debug, and extend: adding
a new capability means writing a new skill or tool, not redesigning an agent coordination
protocol.

% -----------------------------------------------------------------------
\subsection{Skills}
\label{subsec:skills}
% -----------------------------------------------------------------------

Skills are specialized instruction sets that the master agent executes inline. They are
not separate agents --- they run inside the master agent's context using the shared tool
set. Each skill provides domain-specific knowledge for one stage of the pipeline, encoding
expert-level HVAC drawing reading conventions and ontology generation rules that a general
purpose language model would not apply consistently on its own.

\begin{table}[htbp]
  \caption{Master agent skills and their roles in the pipeline.}
  \label{tab:skills}
  \centering
  \begin{tabular}{ll}
    \toprule
    \textbf{Skill} & \textbf{Purpose} \\
    \midrule
    skill-ductwork            & Identify and place all ductwork from HVAC drawings as the spatial foundation \\
    skill-hvac-equipments     & Place HVAC equipment (fans, coils, dampers, sensors) onto the established duct system \\
    skill-bacnet-points       & Extract BACnet point data from CSV exports and map points to grid equipment \\
    skill-ontology-generation & Generate ASHRAE 223P Python code from the populated grid state \\
    skill-ontology-validation & Iteratively fix and validate ontology code to produce a clean TTL file \\
    skill-read-code           & Load error-resolution lessons from past coding sessions to avoid repeating mistakes \\
    \bottomrule
  \end{tabular}
\end{table}

The skills follow a strict dependency chain that mirrors the physical reality of HVAC
commissioning. Ductwork is placed first because all equipment must be anchored to a duct
segment --- the duct grid is the spatial foundation. Equipment placement follows, using
the duct coordinates established in the previous step. BACnet mapping comes next, attaching
live point metadata to the equipment already on the grid. Ontology generation reads the
fully populated grid and produces Python code that models the system in ASHRAE 223P.
Ontology validation then enters an iterative fix loop, executing and correcting the Python
code until it produces a clean TTL file with no runtime errors. Each skill assumes the
previous one has completed successfully; the master agent is responsible for detecting
failures and deciding whether to retry or report the issue.

% -----------------------------------------------------------------------
\subsection{Tools}
% -----------------------------------------------------------------------

The master agent has access to approximately 25 tools organized into ten categories.
Tools handle all interactions with external systems, data storage, and the agent's own
operational state --- the skills provide the reasoning; the tools provide the execution.

\begin{table}[htbp]
  \caption{Master agent tools grouped by category.}
  \label{tab:tools}
  \centering
  \small
  \begin{tabular}{p{3.2cm}p{4.5cm}p{6.5cm}}
    \toprule
    \textbf{Category} & \textbf{Tools} & \textbf{Description} \\
    \midrule
    Grid Sync
      & \texttt{sync\_graphivac\_to\_agent},
        \texttt{sync\_agent\_to\_graphivac}
      & Bidirectional sync between agent state and Graphivac canvas \\
    \addlinespace
    Grid CRUD
      & \texttt{add\_component},
        \texttt{add\_components\_batch},
        \texttt{delete\_component},
        \texttt{delete\_components\_batch},
        \texttt{read\_internal\_grid}
      & Create, read, and delete components in the internal grid \\
    \addlinespace
    Metadata
      & \texttt{write\_metadata},
        \texttt{write\_metadata\_batch}
      & Attach BACnet point data to grid equipment \\
    \addlinespace
    Ontology
      & \texttt{search\_class\_mapping},
        \texttt{scan\_python\_files\_filtered},
        \texttt{read\_ontology},
        \texttt{write\_ontology},
        \texttt{execute\_ontology},
        \texttt{read\_prompt}
      & ASHRAE 223P library lookup, code read/write/execute \\
    \addlinespace
    Ontology Signals
      & \texttt{checkpoint\_code},
        \texttt{exit\_generator\_success},
        \texttt{exit\_generator\_failure},
        \texttt{exit\_validator\_success},
        \texttt{exit\_validator\_failure}
      & Snapshot versioning and ontology task completion signaling \\
    \addlinespace
    Neo4j
      & \texttt{load\_ttl\_to\_neo4j}
      & Import validated TTL semantic model into Neo4j graph database \\
    \addlinespace
    Frontend
      & \texttt{capture\_frontend\_state}
      & Screenshot the live canvas for visual self-correction \\
    \addlinespace
    Ingestion
      & \texttt{ingest\_category\_files}
      & Load user-uploaded files (drawings, CSVs, PDFs) into session \\
    \addlinespace
    State / Progress
      & \texttt{save\_agent\_state},
        \texttt{get\_agent\_state},
        \texttt{update\_step},
        \texttt{update\_status}
      & Track processing progress and persist state across turns \\
    \addlinespace
    Loop Control
      & \texttt{exit\_loop\_level\_2}
      & Terminate the master agent session \\
    \bottomrule
  \end{tabular}
\end{table}

% -----------------------------------------------------------------------
\subsection{Key Technology Decisions}
% -----------------------------------------------------------------------

Several technology choices were made deliberately based on lessons from the iterative
development process. Each decision reflects a tradeoff between capability, cost, and
integration complexity.

\textbf{Google ADK.}
Google's Agent Development Kit was chosen as the agent runtime for its native Gemini
integration and its built-in support for agent lifecycle management, tool registration,
session state, and artifact handling. Earlier iterations used LangChain, which imposed
commercial licensing costs for production use; LangGraph was evaluated but abandoned for
the same reason. ADK provides the necessary infrastructure at lower total cost and with
tighter coupling to the Gemini model family.

\textbf{CopilotKit.}
CopilotKit implements the Agent-to-UI protocol, enabling the Python ADK backend to stream
agent thoughts and actions to the Next.js frontend in real time. Without this library,
building the real-time HUD would have required custom WebSocket infrastructure, additional
message serialization, and a stateful server --- a significant engineering effort with no
research value. CopilotKit handles this concern completely, allowing development to focus
on the agent pipeline itself.

\textbf{Neo4j with Neosemantics (n10s).}
ASHRAE 223P is an RDF ontology; its natural storage format is a graph database with
native RDF import capabilities. Neo4j with the Neosemantics (n10s) plugin satisfies both
requirements: it stores the TTL file as a native graph, and it provides Cypher --- a
mature, expressive query language --- for graph traversal and analysis. Alternative
approaches such as storing TTL as flat files or using a triple store were considered, but
Neo4j's combination of graph-native storage and strong ecosystem tooling made it the
pragmatic choice.

\textbf{Sigma.js for Graph Visualization.}
Sigma.js renders large graphs using WebGL, providing smooth performance on semantic models
with hundreds of nodes and relationships. Heavier alternatives such as D3 or Gephi were
considered, but both introduce unnecessary complexity for a browser-native visualization
requirement. Sigma.js integrates cleanly with the Next.js frontend and requires no
additional server-side rendering accommodations beyond a standard dynamic import guard.
